\begin{ex}
Tính: $356 + 278$.
\shortans{634}
\loigiai{
Ta có $356 + 278 = 634$.
}
\end{ex}

\begin{ex}
Tính: $634 - 356$.
\shortans{278}
\loigiai{
Ta có $634 - 356 = 278$.
}
\end{ex}

\begin{ex}
Tính: $634 - 278$.
\shortans{356}
\loigiai{
Ta có $634 - 278 = 356$.
}
\end{ex}

\begin{ex}
Tính: $485 + 367$.
\shortans{852}
\loigiai{
Ta có $485 + 367 = 852$.
}
\end{ex}

\begin{ex}
Tính: $852 - 485$.
\shortans{367}
\loigiai{
Ta có $852 - 485 = 367$.
}
\end{ex}

\begin{ex}
Tính: $500 - 186$.
\shortans{314}
\loigiai{
Ta có $500 - 186 = 314$.
}
\end{ex}

\begin{ex}
Tính: $1000 - 548$.
\shortans{452}
\loigiai{
Ta có $1000 - 548 = 452$.
}
\end{ex}

\begin{ex}
Tính giá trị biểu thức: $345 + 289 - 176$.
\shortans{458}
\loigiai{
Ta có $345 + 289 = 634$, rồi $634 - 176 = 458$.
}
\end{ex}

\begin{ex}
Tính giá trị biểu thức: $725 - 368 + 247$.
\shortans{604}
\loigiai{
Ta có $725 - 368 = 357$, rồi $357 + 247 = 604$.
}
\end{ex}

\begin{ex}
Tính giá trị biểu thức: $800 - 245 - 155$.
\shortans{400}
\loigiai{
Ta có $800 - 245 = 555$, rồi $555 - 155 = 400$.
}
\end{ex}

\begin{ex}
Điền số thích hợp: $275 + \dots = 612$.
\shortans{337}
\loigiai{
Số cần điền: $612 - 275 = 337$.
}
\end{ex}

\begin{ex}
Điền số thích hợp: $804 - \dots = 358$.
\shortans{446}
\loigiai{
Số cần điền: $804 - 358 = 446$.
}
\end{ex}

\begin{ex}
Điền số thích hợp: $\dots + 386 = 720$.
\shortans{334}
\loigiai{
Số cần điền: $720 - 386 = 334$.
}
\end{ex}

\begin{ex}
Điền số thích hợp: $\dots - 452 = 389$.
\shortans{841}
\loigiai{
Số cần điền: $389 + 452 = 841$.
}
\end{ex}

\begin{ex}
Lớp 2A quyên góp được 275 quyển vở, lớp 2B quyên góp được 318 quyển vở. Cả hai lớp quyên góp được bao nhiêu quyển vở?
\shortans{593}
\loigiai{
Cả hai lớp quyên góp: $275 + 318 = 593$ quyển vở.
}
\end{ex}

\begin{ex}
Kho gạo có 950 bao gạo, đã xuất bán 475 bao. Kho còn lại bao nhiêu bao gạo?
\shortans{475}
\loigiai{
Còn lại: $950 - 475 = 475$ bao gạo.
}
\end{ex}

\begin{ex}
Trang trại thu hoạch được 620 quả dưa hấu, đã bán 385 quả. Còn lại bao nhiêu quả dưa hấu?
\shortans{235}
\loigiai{
Còn lại: $620 - 385 = 235$ quả dưa hấu.
}
\end{ex}

\begin{ex}
Hai số có tổng bằng 762, số hạng thứ nhất là 385. Số hạng thứ hai là bao nhiêu?
\shortans{377}
\loigiai{
Số hạng thứ hai: $762 - 385 = 377$.
}
\end{ex}

\begin{ex}
Hai số có hiệu bằng 268, số trừ là 354. Số bị trừ là bao nhiêu?
\shortans{622}
\loigiai{
Số bị trừ: $268 + 354 = 622$.
}
\end{ex}

\begin{ex}
Hai số có hiệu bằng 345, số bị trừ là 812. Số trừ là bao nhiêu?
\shortans{467}
\loigiai{
Số trừ: $812 - 345 = 467$.
}
\end{ex}

\begin{ex}
Một xưởng sản xuất tháng thứ nhất được 450 sản phẩm, tháng thứ hai nhiều hơn tháng thứ nhất 75 sản phẩm. Tháng thứ hai sản xuất được bao nhiêu sản phẩm?
\shortans{525}
\loigiai{
Tháng 2: $450 + 75 = 525$ sản phẩm.
}
\end{ex}

\begin{ex}
Cả hai tháng xưởng sản xuất được bao nhiêu sản phẩm?
\shortans{975}
\loigiai{
Cả hai tháng: $450 + 525 = 975$ sản phẩm.
}
\end{ex}

\begin{ex}
Một tàu chở hàng chở 650 tấn hàng, sau khi bốc dỡ ở cảng A thì trên tàu còn 280 tấn hàng. Đã bốc dỡ ở cảng A bao nhiêu tấn hàng?
\shortans{370}
\loigiai{
Đã bốc dỡ: $650 - 280 = 370$ tấn.
}
\end{ex}

\begin{ex}
Tính: $1000 - 300 - 400 + 150$.
\shortans{450}
\loigiai{
Ta có $1000 - 300 - 400 = 300$, rồi $300 + 150 = 450$.
}
\end{ex}

\begin{ex}
Số lớn nhất có ba chữ số trừ đi số bé nhất có ba chữ số bằng bao nhiêu?
\shortans{899}
\loigiai{
Số lớn nhất là 999, bé nhất là 100. Hiệu: $999 - 100 = 899$.
}
\end{ex}
